\section{The coefficientwise actual-direction bridge}
\label{sec:tpc52-actual-bridge}

We now combine the sampled overlap theorem with the provenance-fixed Mellin
representation.  The conclusion concerns one distinguished coefficient
density and the coefficientwise pre-terminal Hilbert space.  It is not a
lower-frame theorem for arbitrary signed Mellin densities.

\subsection{The exact actual-to-profile dictionary}

The sampled theorem of \cref{sec:tpc52-profile-overlap} is abstract until its
weights and profiles are identified coordinate by coordinate with
\eqref{eq:tpc52-exact-physical-energy}.  We therefore isolate, rather than
hide, the required application hypothesis.  For notational economy it is
stated on one fixed row-residue group; its finitely many orbit-residue cells
form \(\cA_0\).  A fixed finite collection of row-residue groups is treated
by orthogonal summation.

\begin{assumption}[Exact actual-profile identification]
\label{ass:tpc52-exact-profile-dictionary}
The following data agree on every retained coordinate
\(i=(a,m,r,j,\gamma)\).
\begin{enumerate}[label=\textup{(I\arabic*)},leftmargin=3.2em]
 \item The opposite-row universe is one complete positive-density
 admissible prime--squarefree residue cell in the nondegenerate interior
 regime of TPC--50, \(m=\ell d\) is injective there, and no independent
 opposite-row pruning has occurred.  The only deletion is
 \(M_X(m,\gamma)=K_{m\gamma}\).
 \item The base weight is
 \(\omega_{\gamma,X}=(\log\ell_\gamma)^{2e}\),
 \(e\in\{0,1\}\).  Every fixed dyadic, unary, and residue amplitude not
 contained in the outer coefficient is placed in the named functions
 \(B_{\beta_m,a}\) and \(F_{\beta_m,a}\).
 \item There are outer source coefficients \(a_m\), independent of the
 retained cofactor label \(r\), and diagonal carrier
 coordinates \(c_{X,a}(i)\) such that
 \begin{align}
  |c_{X,a}(i)|^2
  &={}
  M_X(m,\gamma)|a_m|^2\omega_{\gamma,X}
  |B_{\beta_m,a}(q_\gamma,j/J)|^2,
  \label{eq:tpc52-carrier-coordinate-dictionary}\\
  |\gamma_m^{(1)}|^2K_{m\gamma}
  |\Xi_{m,\gamma}(j)|^2
  |\mathcal W_{m,\gamma}(j/J)|^2
  &={}
  M_X(m,\gamma)|a_m|^2\omega_{\gamma,X}
  |F_{\beta_m,a}(q_\gamma,j/J)|^2.
  \label{eq:tpc52-profile-coordinate-dictionary}
 \end{align}
 The direct cells are disjoint and these identities use the same physical
 atomic normalization as \eqref{eq:tpc52-coefficient-space}.
\end{enumerate}
\end{assumption}

This assumption is a checkable dictionary, not a conclusion of the cited
papers.  TPC--25 supplies the constituent Mellin factorizations and
TPC--50/51 supply the complete-row sampling and mask estimates on their
declared cells; neither reference proves that an arbitrary inherited packet
satisfies all three items above.  If the dictionary is not verified, the
results below remain conditional analytic criteria for that packet.

\subsection{Carrier energy and the diagonal layer norm}

Let \(\cU_X\) be the finite set of retained source/cofactor labels
\(u=(m,r)\) in the fixed row-residue group.  For each \(u=(m,r)\), put
\(a_u:=a_m\), where \(a_m\) contains the literal source coefficient and
every factor independent of the opposite row and orbit coordinates.  Thus
the same source amplitude is repeated on every retained cofactor fiber, as
in \eqref{eq:tpc52-actual-vector-profile}.  The opposite-row base weight is
\(\omega_{\gamma,X}\), and the literal mask is
\(M_X(m,\gamma)\).  With the notation of
\cref{sec:tpc52-profile-overlap}, define the unnormalized carrier energy
\begin{equation}
 \mathscr C_X^L
 :=E_XJ\sum_{u=(m,r)\in\cU_X}
       |a_u|^2\mathscr B_X^M(m).
 \label{eq:tpc52-actual-carrier-energy}
\end{equation}
The physical profile energy is
\begin{equation}
 \mathscr P_X^L
 :=E_XJ\sum_{u=(m,r)\in\cU_X}
       |a_u|^2\mathscr P_X^M(m).
 \label{eq:tpc52-actual-profile-energy}
\end{equation}
Here the same source row may occur with any number of cofactors.  No
uniformity of that multiplicity is assumed.

Whenever \(\mathscr C_X^L>0\), define the aggregate actual overlap
\begin{equation}
 \Theta_X^{\rm act}
 :=\frac{\mathscr P_X^L}{\mathscr C_X^L}.
 \label{eq:tpc52-aggregate-actual-overlap}
\end{equation}
On a single direct cell this is the sampled instantiation of
\(\Theta(\mathbf W;c)\) from
\eqref{eq:tpc52-profile-overlap-ratio}; for a fixed finite cell family it is
the corresponding carrier-energy-weighted aggregate.

Summing \eqref{eq:tpc52-profile-coordinate-dictionary} over the disjoint
physical coordinates, and using \(K_{m\gamma}^2=K_{m\gamma}\), gives
\begin{equation}
 \boxed{
 \mathscr P_X^L
 =\|G_X^L\|_{\cH_X^{\rm coeff}}^2
 =\|\cS_X^L\one\|_{\cH_X^{\rm coeff}}^2.}
 \label{eq:tpc52-Q-equals-physical-energy}
\end{equation}
No source sum, prime sum, or residual grouping occurs in
\eqref{eq:tpc52-Q-equals-physical-energy}.  In particular, this equality is
licensed by the displayed coordinate dictionary, not by a change of
normalization after the fact.

The point of real-axis Mellin inversion is that its scale factors have unit
modulus.  More precisely, after the absolute values of the fixed Mellin
transform densities have been placed in \(d\mu_X\), every continuous layer
on a direct cell has the form
\begin{equation}
 A_{X,\theta}^L(i)
 =c_{X,a}(i)\exp(i\Phi_{X,a,\theta}(i)),
 \qquad |\exp(i\Phi_{X,a,\theta}(i))|=1,
 \label{eq:tpc52-unimodular-layer-form}
\end{equation}
on its finite coordinate set.  The vector \(c_{X,a}\) includes the literal
mask and the common compact orbit carrier, but not the physical factors that
are reconstructed by Mellin inversion.  Therefore
\begin{equation}
 \|A_{X,\theta}^L\|=R_{X,a}
 \label{eq:tpc52-layer-carrier-norm}
\end{equation}
is independent of the continuous Mellin parameter on that cell, and
\begin{equation}
 \mathscr C_X^L=\sum_{a\in\cA_0}R_{X,a}^2.
 \label{eq:tpc52-carrier-cell-square}
\end{equation}

For each nonzero cell put
\begin{equation}
 M_{a,B}:=\int_{\Theta_a}w_B(\theta)\,d\mu_{X,a}(\theta).
 \label{eq:tpc52-cell-weighted-Mellin-mass}
\end{equation}
The scale parameters occur only in the phases in
\eqref{eq:tpc52-unimodular-layer-form}.  The transform densities come from a
fixed finite library of nonzero smooth shapes; more generally their weighted
total-variation masses depend continuously on normalized parameters in a
fixed compact set.  Schwartz decay gives the upper bound below, while
nonzero mass and compactness give the lower bound.  Thus there are constants
\begin{equation}
 0<m_B\le M_{a,B}\le M_B<\infty
 \label{eq:tpc52-fixed-Mellin-mass-bounds}
\end{equation}
for all active cells and every \(X\).  A zero-measure template is discarded
before the active-cell family is declared.

\begin{proposition}[Finite-cell carrier identity]
\label{prop:tpc52-finite-cell-carrier-identity}
For the provenance-fixed direction,
\begin{equation}
 \cE_{1,B}^L(\one)=
 \sum_{a\in\cA_0}M_{a,B}R_{X,a}.
 \label{eq:tpc52-envelope-cell-sum}
\end{equation}
Consequently there are fixed constants \(0<c_B\le C_B<\infty\) such
that
\begin{equation}
 c_B\mathscr C_X^L
 \le \bigl(\cE_{1,B}^L(\one)\bigr)^2
 \le C_B\mathscr C_X^L.
 \label{eq:tpc52-envelope-carrier-equivalence}
\end{equation}
\end{proposition}

\begin{proof}
Equation \eqref{eq:tpc52-layer-carrier-norm} makes the norm in the definition
of \(\cE_{1,B}^L\) independent of the continuous Mellin parameter.  Integrate
on each direct cell and then sum over the finite disjoint union to obtain
\eqref{eq:tpc52-envelope-cell-sum}.  The lower bound follows from
\[
 \left(\sum_aM_{a,B}R_{X,a}\right)^2
 \ge m_B^2\sum_aR_{X,a}^2.
\]
For the upper bound, Cauchy--Schwarz in the fixed cell index gives
\[
 \left(\sum_aM_{a,B}R_{X,a}\right)^2
 \le\left(\sum_aM_{a,B}^2\right)
       \left(\sum_aR_{X,a}^2\right).
\]
Use \eqref{eq:tpc52-fixed-Mellin-mass-bounds} and
\eqref{eq:tpc52-carrier-cell-square}.
\end{proof}

This proposition identifies exactly what Schwartz decay, finite physical
support, and Bochner integrability accomplish.  They make the weighted
Mellin mass finite and compare the outer \(L^1(d\mu_X)\) envelope with the
\emph{layer carrier energy}.  They do not compare that carrier with the
signed reassembled physical output.  The latter comparison is precisely the
profile overlap proved in the preceding section.

\subsection{The uniformly active actual direction}

\begin{theorem}[Conditional coefficientwise actual Mellin bridge]
\label{thm:tpc52-coefficientwise-actual-bridge}
Assume the inherited coefficient interface of
\cref{sec:tpc52-interface},
Assumption~\ref{ass:tpc52-exact-profile-dictionary}, the row-sampling and
literal-mask hypotheses
\textup{(S1)}--\textup{(S3)}, the finite direct-cell diagonal form
\eqref{eq:tpc52-unimodular-layer-form}, and
Assumption~\ref{ass:tpc52-profile-nondegeneracy}, and suppose
\(\mathscr C_X^L>0\).  Then for all sufficiently large \(X\),
\(\Theta_X^{\rm act}\ge\rho_*/2\),
\begin{equation}
 \|\cS_X^L\one\|_{\cH_X^{\rm coeff}}^2
 \ge \frac{\rho_*}{2}\mathscr C_X^L,
 \label{eq:tpc52-actual-output-lower}
\end{equation}
and
\begin{equation}
 \boxed{
 \kappa_{\rm can,B}^L(X)
 =\frac{\bigl(\cE_{1,B}^L(\one)\bigr)^2}
   {\|\cS_X^L\one\|_{\cH_X^{\rm coeff}}^2}
 \le \frac{2C_B}{\rho_*}=O_{B,H_0,\cA_0,\mathbf W}(1).}
 \label{eq:tpc52-actual-condition-bound}
\end{equation}
The conclusion is uniform in arbitrary concentration of the actual source
coefficients and arbitrary finite multiplicity of the retained cofactors
\(r\).  The symmetric right coefficient has the same conclusion; a fixed
orthogonal left--right sum changes only the displayed constant.
\end{theorem}

\begin{proof}
Use \eqref{eq:tpc52-Q-equals-physical-energy} and apply
\cref{cor:tpc52-arbitrary-source-r-concentration} with the actual outer
family \(a_{(m,r)}:=a_m\).  This proves
\eqref{eq:tpc52-actual-output-lower}.  The upper inequality in
\eqref{eq:tpc52-envelope-carrier-equivalence} then gives
\eqref{eq:tpc52-actual-condition-bound}.  The proof is fiberwise before the
outer direct sum, so neither source concentration nor cofactor multiplicity
introduces a counting loss.
\end{proof}

The theorem controls the one density fixed by provenance.  It does not
contradict the finite-rank kernel theorem for arbitrary continuous Mellin
densities, and it does not produce a uniform lower frame on the entire
growing translate family.  The orbit, source, opposite-row, and cofactor
labels prevent the artificial scalar cross terms responsible for that false
frame problem.

\subsection{A good/bad actual bridge}

The threshold theorem has an immediate actual-direction form.  Let
\(\mathscr C_{X,\rm good}^L\) be the part of
\eqref{eq:tpc52-actual-carrier-energy} on source fibers in
\(\mathfrak B_{\rm good}(\eta_X)\), and set
\begin{equation}
 \sigma_X^L:=\frac{\mathscr C_{X,\rm good}^L}{\mathscr C_X^L}
 \quad(\mathscr C_X^L>0).
 \label{eq:tpc52-actual-good-fraction}
\end{equation}

\begin{corollary}[Energy-selected actual bridge]
\label{cor:tpc52-energy-selected-actual-bridge}
Retain the dictionary, finite-cell, sampling, mask, and profile hypotheses of
\cref{thm:tpc52-coefficientwise-actual-bridge}.  Assume also
\(\mathscr C_X^L>0\), \(\sigma_X^L>0\), and
\(\Delta_X^{\rm samp}=o(\eta_X)\).  Then
\begin{equation}
 \|\cS_X^L\one\|^2
 \ge \frac{\eta_X\sigma_X^L}{2}\mathscr C_X^L
 \label{eq:tpc52-selected-output-lower}
\end{equation}
and
\begin{equation}
 \kappa_{\rm can,B}^L(X)
 \le \frac{2C_B}{\eta_X\sigma_X^L}.
 \label{eq:tpc52-selected-condition-bound}
\end{equation}
In particular, if
\(\eta_X=X^{-\lambda+o(1)}\) and
\(\sigma_X^L=X^{-\varsigma+o(1)}\), then
\begin{equation}
 \kappa_{\rm can,B}^L(X)
 \le X^{\lambda+\varsigma+o(1)}.
 \label{eq:tpc52-selected-condition-exponent}
\end{equation}
\end{corollary}

\begin{proof}
Equation \eqref{eq:tpc52-selected-output-lower} is
\cref{prop:tpc52-threshold-overlap} applied to the actual coefficients.
Combine it with \eqref{eq:tpc52-envelope-carrier-equivalence}; the exponent
form is immediate.
\end{proof}

The factor \(\sigma_X^L\) cannot be omitted.  If the actual source and
cofactor mass is allowed to concentrate entirely on grazing fibers, a good
set of large row density says nothing about the distinguished actual
direction.  Conversely, no equidistribution of the source is needed once a
positive good-carrier fraction is established by a separate energy
selection.

\subsection{Endpoint cost}

At the inherited endpoint,
\begin{equation}
 Q=X^{267/400+o(1)},
 \qquad J=X^{133/400+o(1)},
 \qquad \kappa_0=\frac1{400}.
 \label{eq:tpc52-endpoint-scales}
\end{equation}
For a fixed collection of direct residue cells, fixed smooth shapes, and a
fixed derivative order \(B\), the weighted Mellin masses in
\eqref{eq:tpc52-fixed-Mellin-mass-bounds} have zero polynomial cost.  Orbit
sampling contributes \(O(J^{-1})=X^{-133/400+o(1)}\), and the inherited
literal-mask defect tends to zero.  Hence a fixed
\(\rho_*>0\) gives \eqref{eq:tpc52-actual-condition-bound} with zero positive
polynomial loss and spends none of the \(1/400\) allowance.

For a shrinking good threshold and a shrinking good-carrier fraction, the
exact cost is instead
\begin{equation}
 \lambda+\varsigma
 \label{eq:tpc52-actual-endpoint-cost}
\end{equation}
in the fixed-cell situation.  More generally one adds the transform and
cell-count cost \(2\tau+\varrho\) from
\eqref{eq:tpc52-cell-mass-exponent}.  Therefore the sufficient ledger is
exactly \eqref{eq:tpc52-one-over-four-hundred-ledger}.  This exponent ledger
does not verify its own relative sampling premise: a polynomially shrinking
threshold still requires \(\Delta_X^{\rm samp}=o(\eta_X)\).  A strict margin
\(\lambda<1/400\) makes only the inherited mask-deletion term negligible.
The available fixed-modulus row-sampling error
\(e^{-c\sqrt{\log L}}\) is not \(o(X^{-\lambda})\) for any fixed
\(\lambda>0\).  Thus the polynomial threshold branch remains conditional on
a sharper relative sampling theorem (or it must be replaced by fixed uniform
overlap).

\subsection{Why the TPC--48 profile does not follow}

The present result is deliberately placed before the terminal prime map.  In
TPC--48, a fixed projective layer contains
\begin{equation}
 B_{\omega,m}(\alpha)
 =\sum_{p,r}A_\omega(p,r,m)\e(-pr\alpha)
  v_{\omega,p,r,m},
 \label{eq:tpc52-downstream-row-vector}
\end{equation}
and its displayed envelope has the order
\begin{equation}
 \mathfrak S_{\mathscr R}(\alpha)
 =\int C_{\rm LS}(\phi_\omega,H_0)^{1/2}
 \left(\sum_m\|B_{\omega,m}(\alpha)\|^2\right)^{1/2}
 d|\nu|(\omega).
 \label{eq:tpc52-full-profile-reminder}
\end{equation}
This is not \(\cE_{1,B}^L(\one)\).  Four operations intervene:
\begin{enumerate}[label=\textup{(D\arabic*)},leftmargin=3.2em]
 \item the prime variable is summed coherently at fixed \((m,r)\) before
       the Hilbert norm is taken;
 \item the outer total-variation integral is taken after a nonlinear square
       root of the layer energy;
 \item \(\omega\) contains projective-mask and terminal parameters that are
       not direct physical Mellin cells; and
 \item \(C_{\rm LS}^{1/2}\) is a profile seminorm weight, not the norm of a
       canonical Mellin coefficient.
\end{enumerate}
The diagonal identity survives unitary coordinate phases, but it does not
survive an arbitrary coherent map with a kernel.  Thus
\cref{thm:tpc52-coefficientwise-actual-bridge} cannot be moved through
\eqref{eq:tpc52-downstream-row-vector} by Minkowski in reverse.

After the present bridge, the remaining interfaces are, in dependency order:
\begin{enumerate}[label=\textup{(O\arabic*)},leftmargin=3.2em]
 \item a lower or two-sided coherent-prime square at fixed \((m,r)\), stable
       under the actual Mellin carrier projection;
 \item control of the many-to-one map grouping \(s=d(m)r\), including its
       fiber multiplicity and possible cancellation;
 \item localization of the completed frozen shift family back to the one
       physical intercept \(h_0\); and
 \item identification of the residual low-sign Gram with the native affine
       physical Gram without discarding its cross terms.
\end{enumerate}
None is proved by profile nondegeneracy.  In particular, the present theorem
contains no fixed-shift prime-pair lower bound, no parity-breaking estimate,
and no twin-prime conclusion.
